%============================================================
% CHƯƠNG 7: ĐẠO LÃNH ĐẠO
%============================================================
\chapter{Đạo Lãnh Đạo (The Way of Leadership)}

\begin{center}
  \textit{\color{truyenblue}``The greatest leader is not necessarily the one who does the greatest things. He is the one who gets the people to do the greatest things.''}\\
  \textit{(Nhà lãnh đạo vĩ đại nhất không nhất thiết là người làm nên những điều vĩ đại, mà là người khiến mọi người cùng làm nên những điều vĩ đại.)}\\[4pt]
  \small\color{truyengold}--- Ronald Reagan ---
\end{center}
\ngancach

\section{Ăn Sau Cùng}
\begin{truyen}{Ăn Sau Cùng}{Phục Vụ}
\chuhoa{T}{ }rong một trại lính thời chiến, sĩ quan chỉ huy luôn là người ăn sau cùng.\\
\textit{(In a wartime military camp, the commanding officer was always the last to eat.)}

Lính mới hỏi: ``Tại sao thủ trưởng không ăn trước cho nóng?''\\
\textit{(A new soldier asked: ``Why does the commander not eat first while the food is hot?'')}

``\textbf{Vì tao ăn lương của họ. Lương thủ trưởng cao hơn là để tao chịu trách nhiệm nhiều hơn, không phải để tao hưởng thụ nhiều hơn}.''\\
\textit{(``\textbf{Because I am paid by them. The commander earns more to bear more responsibility, not to enjoy more}.'')}

Từ hôm đó, toàn đội sẵn sàng đi vào bất kỳ trận đánh nào vị sĩ quan chỉ huy dẫn đầu.\\
\textit{(From that day on, the whole unit was willing to follow the commander into any battle.)}
\end{truyen}
\begin{baihoc}
  Lãnh đạo thực sự được đo bằng mức độ phục vụ của họ cho người dưới quyền, không phải bằng đặc quyền họ tận hưởng. Người đặt lợi ích tập thể lên trên cá nhân mình sẽ được đội ngũ tuyệt đối tin tưởng.
\end{baihoc}

\section{Lỗi Của Sếp}
\begin{truyen}{Lỗi Của Sếp}{Trách Nhiệm}
\chuhoa{T}{ }rước khách hàng quan trọng, nhân viên mắc sai lầm lớn làm hỏng toàn bộ buổi thuyết trình.\\
\textit{(In front of an important client, an employee made a major mistake that ruined the entire presentation.)}

Giám đốc đứng dậy: ``Lỗi này là của tôi. Tôi đã không kiểm duyệt tài liệu kỹ trước khi đến đây.''\\
\textit{(The director stood up: ``This mistake is mine. I did not review the materials thoroughly before coming here.'')}

Sau khi khách về, nhân viên nói: ``Thưa sếp, thực ra là lỗi của em.''\\
\textit{(After the client left, the employee said: ``Boss, it was actually my fault.'')}

Sếp gật đầu: ``Tôi biết. Nhưng trước mặt khách hàng, \textbf{lỗi của đội là lỗi của người lãnh đạo. Sau cửa này, chúng ta mới nói chuyện về việc cải thiện}.''\\
\textit{(The boss nodded: ``I know. But in front of clients, \textbf{the team's mistake is the leader's mistake. Behind this door, we talk about improvement}.'')}
\end{truyen}
\begin{baihoc}
  Lãnh đạo giỏi bảo vệ đội ngũ trước ngoại lực và phản hồi nội bộ sau cửa đóng. Đó là cách xây dựng lòng trung thành không thể mua bằng tiền lương.
\end{baihoc}

\section{Cuộc Họp Không Có Ghế Đầu}
\begin{truyen}{Cuộc Họp Không Ghế Đầu}{Bình Đẳng}
\chuhoa{V}{ }ị giám đốc mới bước vào phòng họp và đặt ghế của mình vào giữa vòng tròn thay vì đầu bàn.\\
\textit{(The new director entered the meeting room and placed his chair in the middle of a circle rather than at the head of the table.)}

Nhân viên ngơ ngác nhìn nhau.\\
\textit{(Employees looked at each other in confusion.)}

``Ý kiến hay nhất không phải lúc nào cũng đến từ người ngồi đầu bàn,'' ông giải thích. ``\textbf{Tôi muốn nghe từ mọi phía, không chỉ từ người gần ghế tôi nhất}.''\\
\textit{(``The best ideas do not always come from whoever sits at the head,'' he explained. ``\textbf{I want to hear from every direction, not only from those closest to my chair}.'')}

Buổi họp đó tạo ra nhiều ý tưởng đột phá nhất trong lịch sử công ty.\\
\textit{(That meeting generated the most groundbreaking ideas in the company's history.)}
\end{truyen}
\begin{baihoc}
  Lãnh đạo biết lắng nghe tạo ra môi trường nơi tài năng thực sự được phát huy. Chỗ ngồi không tạo nên quyền lực thực sự, mà chính khả năng khơi dậy tiềm năng của người khác mới là nguồn quyền lực bền vững nhất.
\end{baihoc}

\section{Người Mang Dù}
\begin{truyen}{Người Mang Dù}{Tầm Nhìn}
\chuhoa{T}{ }rong đoàn quân hành quân dưới mưa, vị tướng mang theo chiếc dù nhưng không che cho mình.\\
\textit{(In a marching army unit in the rain, the general carried an umbrella but did not hold it over himself.)}

Anh lính cận vệ định mở dù ra che cho tướng quân.\\
\textit{(The guard tried to open the umbrella to shelter the general.)}

``Thôi. Ba nghìn người đang ướt vì tao dẫn đường sai. \textbf{Tao không thể khô ráo khi họ ướt vì tao}.''\\
\textit{(``Leave it. Three thousand men are wet because I led them the wrong way. \textbf{I cannot stay dry while they are wet because of me}.'')}

Sau cuộc hành quân, toàn quân \textbf{tình nguyện đi theo ông vào trận khó nhất} mà không cần lệnh.\\
\textit{(After the march, the whole army \textbf{volunteered to follow him into the hardest battle} without waiting for orders.)}
\end{truyen}
\begin{baihoc}
  Lãnh đạo chia sẻ gian khổ với đội ngũ không chỉ là đức hạnh mà còn là chiến lược. Người sẵn sàng ướt cùng lính của mình sẽ có những người lính sẵn sàng chết vì mình.
\end{baihoc}

\section{Câu Hỏi Của Người Mới}
\begin{truyen}{Câu Hỏi Của Người Mới}{Khiêm Tốn}
\chuhoa{G}{ }iám đốc cũ kỹ của tập đoàn lớn mời nhân viên mới nhất vào văn phòng.\\
\textit{(The veteran director of a large corporation invited the newest employee to his office.)}

``Em làm chưa đầy một tháng. Em thấy công ty mình có vấn đề gì?''\\
\textit{(``You have been here less than a month. What problems do you see in the company?'')}

Người nhân viên ngập ngừng. ``Thưa anh, em chưa có đủ thông tin để...''\\
\textit{(The employee hesitated. ``Sir, I do not have enough information yet to...'')}

``\textbf{Đó chính xác là lý do tôi mời em vào. Người trong cuộc nhìn lâu sẽ không còn thấy gì nữa. Mắt mới thấy những thứ mắt cũ đã bỏ qua}.''\\
\textit{(``\textbf{That is exactly why I invited you. Those who have been inside too long stop seeing anything. Fresh eyes notice what old eyes have overlooked}.'')}
\end{truyen}
\begin{baihoc}
  Lãnh đạo khôn ngoan không ngại hỏi người cấp dưới hay mới nhất. Sự khiêm tốn trí tuệ là dấu hiệu của người đứng đầu có đủ tự tin để biết mình không biết tất cả.
\end{baihoc}

\section{Quyết Định Lúc Nửa Đêm}
\begin{truyen}{Quyết Định Lúc Nửa Đêm}{Can Đảm}
\chuhoa{L}{ }úc nửa đêm, giám đốc nhận báo cáo: sản phẩm mới có lỗi nguy hiểm dù đã ra thị trường.\\
\textit{(At midnight, the director received a report: the newly launched product had a dangerous defect.)}

Cố vấn pháp lý khuyên im lặng vì chưa có người bị hại và thu hồi sẽ tốn hàng triệu đô.\\
\textit{(Legal counsel advised silence since no one had been harmed yet and a recall would cost millions.)}

Ông ngồi im một lúc rồi ký lệnh thu hồi ngay trong đêm.\\
\textit{(He sat quietly for a moment then signed the recall order that very night.)}

``\textbf{Tiền có thể kiếm lại. Lòng tin của khách hàng mà mất thì không miếng quảng cáo nào mua lại được}.''\\
\textit{(``\textbf{Money can be earned back. But customer trust, once lost, cannot be bought back by any advertisement}.'')}

Quyết định đó làm cổ phiếu giảm mạnh tháng đó nhưng \textbf{thương hiệu tồn tại thêm năm mươi năm}.\\
\textit{(That decision caused the stock to drop sharply that month but \textbf{the brand survived for fifty more years}.)}
\end{truyen}
\begin{baihoc}
  Lãnh đạo thực sự biết rằng đạo đức không phải là đối lập với kinh doanh mà là nền tảng của nó. Quyết định đứng về phía đúng khi tốn kém nhất mới là bằng chứng của nhân cách lãnh đạo.
\end{baihoc}

\section{Người Lau Bảng}
\begin{truyen}{Người Lau Bảng}{Gương Mẫu}
\chuhoa{V}{ }ị giáo sư đại học nổi tiếng, sau mỗi buổi giảng, tự mình lau bảng trước khi ra về.\\
\textit{(The famous university professor, after every lecture, cleaned the board himself before leaving.)}

Sinh viên hỏi: ``Thầy có trợ giảng lo việc đó rồi thưa thầy.''\\
\textit{(A student said: ``There are teaching assistants to do that, professor.'')}

``Tôi biết. Nhưng \textbf{tôi không muốn người đến sau phải nhìn thấy sự bừa bộn của tôi}. Điều đó không liên quan đến chức danh.''\\
\textit{(``I know. But \textbf{I do not want those who come after me to face my mess}. That has nothing to do with title.'')}

Nhiều thế hệ sinh viên sau đó nhớ bài học này hơn bất kỳ bài giảng nào.\\
\textit{(Generations of students remembered this lesson more than any lecture.)}
\end{truyen}
\begin{baihoc}
  Lãnh đạo không phải là người được phục vụ mà là người làm gương. Khi người đứng đầu tự làm những việc nhỏ nhặt, họ dạy đội ngũ nhiều hơn bất kỳ bài phát biểu nào.
\end{baihoc}

\section{Hai Người Xin Việc}
\begin{truyen}{Hai Người Xin Việc}{Đánh Giá Đúng}
\chuhoa{H}{ }ai ứng viên vào vòng cuối cho vị trí quản lý cấp cao: một người bằng cấp xuất sắc, một người thành tích khiêm tốn.\\
\textit{(Two candidates reached the final round for a senior management position: one had excellent qualifications, the other modest credentials.)}

Giám đốc nhân sự gọi cả hai vào, đưa cho mỗi người một cái chổi và nhờ dọn phòng họp đang bừa bộn.\\
\textit{(The HR director called them both in, handed each a broom and asked them to clean the messy meeting room.)}

Người bằng cấp cao khựng lại: ``Thưa anh, đây không phải vị trí của tôi.''\\
\textit{(The highly qualified candidate paused: ``Sir, this is not my position.'')}

Người kia cầm chổi quét ngay.\\
\textit{(The other took the broom and started sweeping.)}

``\textbf{Người lãnh đạo giỏi không từ chối bất kỳ việc gì cần làm. Chúc mừng em được nhận}.''\\
\textit{(``\textbf{A good leader refuses no work that needs doing. Congratulations, you are hired}.'')}
\end{truyen}
\begin{baihoc}
  Nhân cách lãnh đạo không nằm ở bằng cấp mà nằm ở thái độ với công việc. Người sẵn sàng làm bất kỳ điều gì vì tập thể mới xứng đáng dẫn dắt tập thể đó.
\end{baihoc}

\section{Cuộc Chia Tay Của Trưởng Nhóm}
\begin{truyen}{Cuộc Chia Tay Của Trưởng Nhóm}{Nuôi Dưỡng}
\chuhoa{K}{ }hi trưởng nhóm giỏi nhất công ty thông báo nghỉ việc để khởi nghiệp, CEO không cố giữ ông lại.\\
\textit{(When the company's best team leader announced his resignation to start his own business, the CEO did not try to stop him.)}

Thay vào đó, CEO tổ chức tiệc chia tay lớn và viết thư giới thiệu chi tiết cho ông.\\
\textit{(Instead, the CEO threw a big farewell party and wrote a detailed recommendation letter for him.)}

Người ngoài hỏi: ``Sao anh không giữ ông ấy lại?''\\
\textit{(An outsider asked: ``Why did you not keep him?'')}

``\textbf{Tôi thuê ông ấy để phát triển, không phải để đứng yên. Nếu tôi giữ ông ấy lại vì sợ mất, tôi sẽ là người đã phá giới hạn của ông ấy. Một lãnh đạo tốt không nuôi con chim trong lồng}.''\\
\textit{(``\textbf{I hired him to grow, not to stand still. If I held him back out of fear of losing him, I would have been the one limiting him. A good leader does not keep a bird in a cage}.'')}
\end{truyen}
\begin{baihoc}
  Lãnh đạo tốt nhất giúp người của mình lớn lên dù điều đó có nghĩa là họ sẽ rời đi. Sự phát triển của người dưới quyền là di sản thực sự của người đứng đầu.
\end{baihoc}

\section{Ghế Trống}
\begin{truyen}{Ghế Trống}{Gần Gũi}
\chuhoa{V}{ }ị giám đốc có thói quen để một chiếc ghế trống trong mọi cuộc họp cấp cao.\\
\textit{(The director had a habit of leaving one empty chair in every senior meeting.)}

Khi được hỏi, ông giải thích: ``\textbf{Chiếc ghế đó đại diện cho khách hàng. Trước khi ra bất kỳ quyết định nào, tôi tự hỏi người ngồi ở đó sẽ nghĩ gì}.''\\
\textit{(When asked, he explained: ``\textbf{That chair represents the customer. Before any decision, I ask myself what the person sitting there would think}.'')}

Chính nhờ chiếc ghế trống ấy, công ty tránh được nhiều quyết định lấy lợi ích nội bộ làm trọng tâm.\\
\textit{(Thanks to that empty chair, the company avoided many decisions driven purely by internal interests.)}

Nhân viên mới vào thường nhớ mãi hình ảnh \textbf{chiếc ghế không người mà có sức mạnh hơn cả ban giám đốc cộng lại}.\\
\textit{(New employees always remembered the image of \textbf{the unoccupied chair that held more power than the entire board combined}.)}
\end{truyen}
\begin{baihoc}
  Lãnh đạo tốt nhất không bao giờ quên người mà tổ chức của mình phục vụ. Đặt lợi ích của khách hàng hay cộng đồng vào trung tâm mọi quyết định là nền tảng của sự lãnh đạo bền vững.
\end{baihoc}
